\section{Step 6: Check the pole equation in the inverted Rosen--Morse barrier}
\label{app:rosen-morse}

With $\xi=\tanh(\beta x)$, Eq.~\eqref{eq:wkb-schrodinger} for
Eq.~\eqref{eq:rosen-morse-potential} becomes
\begin{equation}
  \frac{\rmd}{\rmd\xi}
  \left[(1-\xi^2)\frac{\rmd\psi}{\rmd\xi}\right]
  +\left[
  s(s+1)
  -\frac{(\rmi k/\beta)^2}{1-\xi^2}
  \right]\psi=0 .
  \label{eq:associated-legendre-rm}
\end{equation}
One solution outgoing at $x\to+\infty$ is
\begin{align}
  \psi_k(x)
  &=(1-\xi^2)^{-\rmi k/(2\beta)}
  \notag\\
  &\quad\times
  {}_2F_1\left(
  -\frac{\rmi k}{\beta}-s,
  -\frac{\rmi k}{\beta}+s+1;
  1-\frac{\rmi k}{\beta};
  \frac{1-\xi}{2}
  \right) .
  \label{eq:rm-hypergeometric}
\end{align}
The hypergeometric connection formula around $\xi=-1$ yields reflected and
incident coefficients.  The incident coefficient is
Eq.~\eqref{eq:rosen-morse-incoming}.  Since the Gamma function has poles at
nonpositive integers, one resonance branch obeys
\begin{equation}
  -\frac{\rmi k}{\beta}+s+1=-n,
  \qquad n=0,1,2,\ldots .
  \label{eq:rm-gamma-condition}
\end{equation}
Substitution of the branch of $s$ fixed by
Eq.~\eqref{eq:rosen-morse-parameters} gives
Eq.~\eqref{eq:rosen-morse-k}.  The other denominator Gamma function produces
the antiresonant branch.

The same Gamma-function ratio is a resummed Voros connection coefficient.
Its poles and zeros encode the infinite periodic Stokes graph without the need
to cross every repeated cell separately.  This is an example of a special
function acting as a closed-form resurgent connection matrix.
